\chapter{TMJ Pain (Jaw Pain)}

\section*{Definition}
Temporomandibular joint (TMJ) pain is discomfort arising from the articulation between the mandibular condyle and the temporal bone, or from the masticatory muscles. It affects approximately 5--12\% of the population, with peak incidence in women aged 20--40 years.

\section*{What Happens in Your Body}
The TMJ is a ginglymoarthrodial joint (hinge and sliding) with an articular disc that distributes load during mastication. Pain may derive from the joint itself (capsulitis, synovitis, disc displacement) or the masticatory muscles (myofascial pain). Parafunctional habits (bruxism, clenching) produce sustained muscle contraction, leading to reduced blood flow, metabolite accumulation, and central sensitization. Disc displacement with reduction produces clicking; without reduction it causes locking and limited mouth opening.

\section*{Causes}
\begin{itemize}
  \item \textbf{Myofascial:} Bruxism (nighttime or daytime clenching), jaw tension, stress, anxiety
  \item \textbf{Joint-related:} Disc displacement (with or without reduction), hypermobility, dislocation, osteoarthritis, rheumatoid arthritis
  \item \textbf{Traumatic:} Mandibular fracture, whiplash, prolonged mouth opening (dental procedures, intubation)
  \item \textbf{Dental:} Malocclusion, missing teeth, poorly fitting dentures, third molar extraction
  \item \textbf{Inflammatory:} Juvenile idiopathic arthritis, gout, pseudogout, psoriatic arthritis
  \item \textbf{Infectious:} Septic TMJ arthritis (rare), odontogenic collection of pus, parotitis
  \item \textbf{Habitual:} Gum chewing, nail biting, lip biting, wide yawning
\end{itemize}

\section*{When to See a Doctor}
See your doctor if:
\begin{itemize}
  \item Jaw pain with limitation of mouth opening (trismus) or locking
  \item Pain with chewing, yawning, or speaking affecting nutrition
  \item Joint clicking, popping, or crepitus with painful locking
  \item Unexplained dental wear or morning headache suggestive of bruxism
\end{itemize}

\section*{Related Symptoms}
\begin{itemize}
  \item Pre-auricular pain and tenderness
  \item Clicking, popping, or crepitus with jaw movement
  \item Limited mouth opening (< 35 mm interincisal distance)
  \item Jaw deviation to one side on opening
  \item Headache, earache, facial pain, and cervical muscle tension
  \item Bruxism-related symptoms: tooth wear, morning headache, masseter hypertrophy
\end{itemize}
\textbf{Important:} TMJ pain can mimic otalgia (earache); a normal otoscopic examination should prompt TMJ evaluation in people with ear pain.

\section*{Self-Care / Home Management}
\begin{itemize}
  \item Soft diet (avoid chewy, crunchy, or tough foods)
  \item Jaw rest with voluntary avoidance of clenching, gum chewing, and wide yawning
  \item Heat or ice packs over the TMJ and masseter muscle for 15 minutes
  \item Over-the-counter NSAIDs (ibuprofen, naproxen) for pain and inflammation
  \item Gentle jaw range-of-motion exercises (without pain)
\end{itemize}

\section*{How Your Doctor Will Evaluate You}
\begin{itemize}
  \item History and physical examination including palpation of TMJ and masticatory muscles
  \item Measurement of maximal interincisal opening and auscultation for joint sounds
  \item Panoramic radiography or cone-beam CT (CBCT) for osseous anatomy and pathology
  \item MRI for disc position, morphology, and joint fluid buildup
  \item Ultrasound for real-time disc dynamics (evolving use)
  \item Psychiatric or sleep evaluation if anxiety, stress, or sleep bruxism is suspected
\end{itemize}

\section*{Treatment Options}
\begin{itemize}
  \item Patient education, jaw rest, and soft diet as first-line
  \item NSAIDs and muscle relaxants (e.g., cyclobenzaprine, diazepam) for acute and subacute pain
  \item Occlusal splint (night guard) for bruxism and disc displacement with reduction
  \item Physical therapy with manual mobilization, posture correction, and exercises
  \item Botulinum toxin injection into masseter and temporalis for refractory bruxism
  \item Arthrocentesis or arthroscopy for joint lavage and lysis of adhesions
  \item Open joint surgery (disc repositioning, discectomy, condyloplasty) for advanced structural disease
\end{itemize}

\section*{References}
\begin{thebibliography}{99}
\bibitem{tmj_pain:uptodate-tmd} Scrivani SJ, et al. Temporomandibular disorders: clinical features and diagnosis. In: UpToDate, Post TW (Ed), UpToDate, Waltham, MA; 2025.
\bibitem{tmj_pain:nejm-tmd} Durham J, et al. ``Temporomandibular Disorders.'' \emph{N Engl J Med}. 2023;389(7):652-662.
\bibitem{tmj_pain:bmj-tmd} Schiffman E, et al. ``Temporomandibular disorders.'' \emph{BMJ}. 2024;385:e079543.
\bibitem{tmj_pain:cochrane-tmd} List T, et al. ``Occlusal splints for TMD.'' \emph{Cochrane Database Syst Rev}. 2023;(8):CD014321.
\bibitem{tmj_pain:jam-dent-tmd} Ohrbach R, et al. ``Diagnostic criteria for TMD.'' \emph{J Am Dent Assoc}. 2024;155(3):234-246.
\bibitem{tmj_pain:harrison-orofacial} Jameson JL, et al. \emph{Harrison\'s Principles of Internal Medicine}. 21st ed. McGraw-Hill; 2022. Chapter 386: Orofacial Pain.
\end{thebibliography}
